\section{Relation to established methods}

The finite-ground construction uses classical congruence closure. Nelson--Oppen and Downey--Sethi--Tarjan established foundational algorithms for ground equality \cite{NelsonOppen1980,DowneySethiTarjan1980}; Kapur related congruence closure to completion \cite{Kapur1997}. The activation filtration records when the equations and represented terms needed by these procedures become available. The locality theorem is the classical subterm property expressed as a depth bound; the Catalan theorem classifies the possible exact profiles under that bound.

Proof production also has a substantial prior literature \cite{NieuwenhuisOliveras2005,NieuwenhuisOliveras2007}. Minimum explanation size and minimum support cardinality differ from maximum term depth \cite{FellnerFontainePaleo2017,FlattEtAl2022}. The two-sided certificate here combines an explanation at the first complete threshold with a finite model below it. Local inference and local theory extensions likewise control the instantiated syntax required for ground reasoning, under their respective structural hypotheses \cite{Ganzinger2001,Sofronie2005}. No extension-locality hypothesis is needed for our finite set of ground equations.

For universal presentations, the action-chain construction links sequential rewrite depth with semigroup derivational space. Olshanskii's invariant counts maximum intermediate word length, and his realization theorems supply the complexity consequences used here \cite{Olshanskii2013}. The skeleton argument accounts for arbitrary one-sorted contexts by the explicit overhead envelope $H$. Example~\ref{ex:chaingap} marks the limit of that correspondence: congruence proofs can reuse established equalities without carrying their entire intermediate syntax inside the current context.

Paper I supplies the algebraic interpretation of the observed interface and the quotient-feedback mechanism \cite{ShcherbakovI}. Here the input terms are arbitrary ground syntax; the connection to incomplete actions is an application. This division lets Paper III cite the query-subterm, first-merge, and threshold-model results directly when describing a concrete certifying implementation.
